Le chapitre le plus admis du programme, et donc celui où la formalisation déplace le plus de
choses.

Comme pour les suites, les deux premiers énoncés vérifient que les définitions de la
bibliothèque sont celles du cours : la limite en un point avec $\varepsilon$ et $\delta$,
la limite en $+\infty$ avec $A$. Le reste du chapitre — opérations sur les limites,
croissances comparées, théorème des valeurs intermédiaires, lien entre le signe de la dérivée
et le sens de variation — est emprunté à la bibliothèque, qui les démontre. Les transcriptions
disent à chaque fois quel argument porte le résultat : la complétude de $\mathbb{R}$ derrière
le théorème des valeurs intermédiaires, les accroissements finis derrière le sens de variation.

Deux points de vigilance. Le taux d'accroissement n'a de sens qu'en $x \ne a$, et c'est le
voisinage épointé, non une convention de calcul, qui l'assure. Et la formule
$(x^n)' = n x^{n-1}$ s'écrit avec un exposant entier naturel : pour $n = 0$, la soustraction
$n - 1$ est tronquée à zéro, ce que le facteur $n$ rend inoffensif.

Deux énoncés du programme ne sont pas formalisés : la réciproque de \og{} dérivable donc
continue \fg{}, dont le contre-exemple est mentionné sans être écrit, et la concavité de la
fonction cube sur les négatifs.
